\documentclass[11pt]{article}
\usepackage[T1]{fontenc}
\usepackage{amsmath,amssymb}
\usepackage{cite}
\usepackage{booktabs}
\usepackage{graphicx}
\usepackage{xurl}
\usepackage{hyperref}
\hypersetup{
  colorlinks=true,
  linkcolor=blue,
  urlcolor=blue,
  citecolor=blue,
  pdfauthor={Jabulani Monchwe, Martin Chanza, Olorato Motitswane, Mfundo Monchwe}
}
\newcommand{\software}{\textsf{ReliableML}}
\newcommand{\code}[1]{\texttt{\small #1}}
\newcommand{\jabulaniauthor}{
\IEEEauthorblockN{Jabulani Monchwe}
\IEEEauthorblockA{\textit{Applied Business Analytics and Decision Making}\\
\textit{North-West University}\\
\textit{IDEA-Research Club Chairperson}\\
Johannesburg, South Africa\\
\href{mailto:monchwej.25@gmail.com}{monchwej.25@gmail.com}}
}
\newcommand{\martinchanzauthor}{
\IEEEauthorblockN{Martin Chanza}
\IEEEauthorblockA{\textit{Applied Business Analytics and Decision Making}\\
\textit{North-West University}\\
Mafikeng, South Africa\\
\href{mailto:martin.chanza@nwu.ac.za}{martin.chanza@nwu.ac.za}}
}
\newcommand{\oloratoauthor}{
\IEEEauthorblockN{Olorato Motitswane}
\IEEEauthorblockA{\textit{Eduvos Midrand Campus -- IT Faculty Data Science, South Africa}\\
\textit{IDEA-Research Club}\\
\href{mailto:oloratomotitswane@gmail.com}{oloratomotitswane@gmail.com}}
}
\newcommand{\mfundoauthor}{
\IEEEauthorblockN{Mfundo Monchwe}
\IEEEauthorblockA{\textit{Eduvos Midrand Campus -- IT Faculty Data Science, South Africa}\\
\textit{IDEA-Research Club}\\
\href{mailto:diditmfundo@gmail.com}{diditmfundo@gmail.com}}
}
\newcommand{\minenhleauthor}{
\IEEEauthorblockN{Minenhle Mchunu}
\IEEEauthorblockA{\textit{Applied Business Analytics and Decision Making}\\
\textit{Eduvos Midrand Campus -- IT Faculty Data Science, South Africa}\\
\textit{IDEA-Research Club}\\
\href{mailto:Mchunu.minenhle02@gmail.com}{Mchunu.minenhle02@gmail.com}}
}
\newcommand{\projecturl}{\url{https://github.com/Jabualani/ML-Reliabilty}}

\title{\software: Methods for Reliable Machine-Learning Workflows}
\author{\contactauthors\\\small MonchweJ.25@gmail.com}
\date{Version 0.1.1 \quad 18 October 2026}
\begin{document}
\maketitle

\begin{abstract}
\software{} is a Python library for attaching data-distribution monitoring,
labeled evaluation, fairness summaries, model-deterioration tracking,
meta-learning, model-agnostic explanations, and statistical reason auditing
to machine-learning workflows. This document summarizes the implemented
methods, their interfaces, and their limits. Its central methodological
qualification is that reliability is not a single scalar: distribution
similarity, predictive performance, calibration, subgroup disparities,
temporal trends, and explanatory stability answer different questions and
must be interpreted separately.
\end{abstract}

\section{Scope and workflow}
Reliable model development begins with appropriate data preparation. Users
may prepare data with ATTP or another tool; \software{} accepts the resulting
features and a user-selected estimator. The library does not claim that
preprocessing, a high validation score, or an explanation by itself guarantees
reliable deployment.

The operating workflow is:
\begin{enumerate}
  \item Fit a model and capture a reference profile of training features.
  \item Predict on incoming features and compare their distributions with the
        reference, optionally saving plots and machine-readable reports.
  \item When labels become available, evaluate predictive quality, calibration,
        statistical associations, and fairness on representative labeled data.
  \item Across dated labeled batches, estimate an accuracy trend.
  \item Where appropriate, use out-of-fold meta-learning, local/global
        explanations, and a holdout reason audit to investigate stability and
        reliance.
\end{enumerate}

\section{Implemented method families}
\begin{description}
  \item[Distribution shift.] Quantile-binned numerical drift statistics,
  categorical homogeneity tests, reference plots, and a labeled
  feature--target divergence diagnostic.
  \item[Evaluation.] Classification metrics including multiclass Matthews
  correlation, confusion counts, feature--target and train--test ANOVA,
  MANOVA and chi-square tests, and an empirical repeated-CV Brier
  bias--variance decomposition.
  \item[Fairness.] Statistical-parity selection-rate, equalized-odds
  true-/false-positive-rate, and predictive-parity positive-predictive-value
  gaps, with group denominators and bootstrap intervals.
  \item[Deterioration.] A linear trend in accuracy across at least three dated,
  labeled batches and a bootstrap interval for the slope.
  \item[Meta-learning.] Cross-fitted base-model meta-features and selectable,
  calibrated, statistical, and neural meta-classifiers. Halving grid and
  randomized search helpers are provided.
  \item[Explainability.] A Monte Carlo permutation estimator of interventional
  Shapley contributions and a LIME-style local weighted linear surrogate.
  \item[Reason audit.] Holdout validity, calibration, perturbation-based
  feature reliance, stability slices, counterfactual sensitivity, logistic
  association, and explanation-alignment diagnostics.
  \item[Pipeline and persistence.] A scikit-learn-style wrapper stores an
  estimator and its fitted monitoring profile, reports prediction-time drift,
  and supports model persistence and exports.
\end{description}

\section{Interpreting evidence}
Each method has different data requirements. Drift monitoring can be
unlabeled, but it does not establish that predictions are wrong. Performance,
calibration, fairness, and deterioration assessment require labels and
appropriate sampling. Counterfactual sensitivity is only meaningful when
replacement values are domain-valid. Explanation agreement measures
consistency, not truth or causality. Logistic odds ratios describe
conditional associations, not causal effects.

Bootstrap intervals quantify sampling variation under their resampling
assumptions. Current intervals generally treat rows as independent; clustered
or time-dependent data may need cluster or block methods. Reported statistics
should be considered alongside data provenance, leakage checks, temporal and
group-aware validation, and domain review. See the standalone method
documents for operational details and assumptions.

\section{Minimal usage}
\begin{verbatim}
from ml_reliability import ReliablePipeline

model = ReliablePipeline(estimator=classifier).fit(X_train, y_train)
predictions, report = model.predict_with_report(X_test)
evaluation = model.evaluate(X_test, y_test)
model.save("reliable-model.pkl")
\end{verbatim}
The same fitted pipeline can be reloaded and used with input data in the
original feature schema. Pickle files must only be loaded from trusted sources.

\section{References}
\begin{thebibliography}{9}
\bibitem{gama}
J. Gama, I. Zliobaite, A. Bifet, M. Pechenizkiy, and A. Bouchachia.
 A survey on concept drift adaptation. \emph{ACM Computing Surveys},
 46(4), 2014. \url{https://doi.org/10.1145/2523813}.
\bibitem{wolpert}
D. H. Wolpert. Stacked generalization. \emph{Neural Networks}, 5(2),
  1992. \url{https://doi.org/10.1016/S0893-6080(05)80023-1}.
\bibitem{hardt}
M. Hardt, E. Price, and N. Srebro. Equality of opportunity in supervised
  learning. \emph{Advances in Neural Information Processing Systems}, 2016.
  \url{https://arxiv.org/abs/1610.02413}.
\bibitem{lime}
M. T. Ribeiro, S. Singh, and C. Guestrin. ``Why should I trust you?''
  Explaining the predictions of any classifier. \emph{Proceedings of KDD},
  2016. \url{https://doi.org/10.1145/2939672.2939778}.
\bibitem{shap}
S. M. Lundberg and S.-I. Lee. A unified approach to interpreting model
  predictions. \emph{Advances in Neural Information Processing Systems}, 2017.
  \url{https://arxiv.org/abs/1705.07874}.
\end{thebibliography}
\end{document}
